\documentclass[10pt,a4paper,ragged2e,withhyper]{altacv}
\geometry{left=1.25cm,right=1.25cm,top=1.5cm,bottom=1.5cm,columnsep=1.2cm,head=14pt,headsep=10pt,footskip=24pt}
\usepackage{paracol}
\usepackage{needspace} % evita títulos huérfanos al final de página (ver Referencias)

% Configuración de fuentes
\iftutex
  \setmainfont{Roboto Slab}
  \setsansfont{Lato}
  \renewcommand{\familydefault}{\sfdefault}
\else
  \usepackage[rm]{roboto}
  \usepackage[defaultsans]{lato}
  \renewcommand{\familydefault}{\sfdefault}
\fi

% Colores
\definecolor{SlateGrey}{HTML}{2E2E2E}
\definecolor{LightGrey}{HTML}{666666}
\definecolor{DarkPastelRed}{HTML}{450808}
\definecolor{PastelRed}{HTML}{8F0D0D}
\definecolor{GoldenEarth}{HTML}{E7D192}
\colorlet{name}{black}
\colorlet{tagline}{PastelRed}
\colorlet{heading}{DarkPastelRed}
\colorlet{headingrule}{GoldenEarth}
\colorlet{subheading}{PastelRed}
\colorlet{accent}{PastelRed}
\colorlet{emphasis}{SlateGrey}
\colorlet{body}{LightGrey}

% Estilos de fuente
\renewcommand{\namefont}{\huge\rmfamily\bfseries}
\renewcommand{\personalinfofont}{\footnotesize}
\renewcommand{\cvsectionfont}{\LARGE\rmfamily\bfseries}
\renewcommand{\cvsubsectionfont}{\large\bfseries}
\renewcommand{\cvItemMarker}{{\small\textbullet}}
\renewcommand{\cvRatingMarker}{\faCircle}

% Configuración para publicaciones
\input{pubs-num.cfg}

% --- Override de \cvtag: respeta el ancho de la columna ---
% El \cvtag original (altacv.cls) es un nodo TikZ de una sola línea sin ancho máximo:
% una aptitud larga se desborda por el borde derecho. Aquí medimos el texto; si su
% ancho natural (+ padding) supera la columna, limitamos el ancho al de la columna y
% dejamos que ENVUELVA dentro del chip. Los tags cortos siguen siendo pills inline.
\newlength{\cvtagnatwidth}
\renewcommand{\cvtag}[1]{%
  \settowidth{\cvtagnatwidth}{#1}%
  \addtolength{\cvtagnatwidth}{2ex}% aprox. inner xsep (1ex por lado)
  \ifdim\cvtagnatwidth>\linewidth
    \tikz[baseline]\node[anchor=base,draw=body!30,rounded corners,%
      inner xsep=1ex,inner ysep=0.75ex,%
      text width=\dimexpr\linewidth-2.5ex\relax,align=left]{#1};%
  \else
    \tikz[baseline]\node[anchor=base,draw=body!30,rounded corners,%
      inner xsep=1ex,inner ysep=0.75ex,text height=1.5ex,text depth=.25ex]{#1};%
  \fi
  % Espacio rompible ENTRE chips (breakpoint + separación visual).
  \penalty0\hspace{2pt plus 2pt}%
}

% \cvtagbox: encierra un grupo de chips en una CAJA del ancho de la columna.
% Sin esto, en una columna de paracol TeX amontona chips en una "línea" más ancha
% que la columna y el conjunto se desborda. El minipage fuerza el quiebre DENTRO
% del ancho disponible → los chips nunca se salen del borde (idea del creator).
\newcommand{\cvtagbox}[1]{%
  \par\noindent\begin{minipage}[t]{\linewidth}\raggedright #1\end{minipage}\par%
}

% Override de \divider: el original no fuerza quiebre de párrafo, y como Jinja
% (trim_blocks) pega `\divider\cvevent{...}` en una línea, el título del siguiente
% cargo se colaba al final de la línea punteada. Aquí la regla va en su PROPIA línea
% (\par antes y después) → el \cvevent siguiente arranca limpio debajo.
\renewcommand{\divider}{%
  \par\smallskip
  {\color{body!30}\hdashrule{\linewidth}{0.6pt}{0.5ex}}%
  \par\smallskip
}

% Override de \cvevent: empuja la UBICACIÓN al margen derecho.
% El original pone fecha y ubicación en dos cajas de 0.5\linewidth, ambas alineadas
% a la izquierda → en ancho completo la ubicación queda al centro con hueco a la
% derecha. Aquí la caja de ubicación se alinea a la derecha ([r]) → fecha a la
% izquierda, ubicación pegada al margen derecho. Solo cambia esa alineación.
\renewcommand{\cvevent}[4]{%
  {\large\color{emphasis}#1\par}%
  \smallskip\normalsize
  \ifstrequal{#2}{}{}{%
    \textbf{\color{accent}#2}\par
    \smallskip}%
  \ifstrequal{#3}{}{}{%
    {\small\makebox[0.5\linewidth][l]%
      {\BeginAccSupp{method=pdfstringdef,ActualText={\datename:}}\cvDateMarker\EndAccSupp{}~#3}%
    }}%
  \ifstrequal{#4}{}{}{%
    {\small\makebox[0.5\linewidth][r]%
      {\BeginAccSupp{method=pdfstringdef,ActualText={\locationname:}}\cvLocationMarker\EndAccSupp{}~#4}%
    }}\par
  \medskip\normalsize
}

% Margen de emergencia: permite a TeX estirar un poco la última pasada antes de
% desbordar, absorbiendo palabras largas de las viñetas (Renal Crónica, etc.).
\setlength{\emergencystretch}{3em}

% --- Marca de agua (header) desde pág. 2 + número de página (centrado) desde pág. 1 ---
% altacv fuerza \pagestyle{empty} en un hook de begin-document; por eso los estilos
% se APLICAN dentro del body (tras \makecvheader), no aquí.
\usepackage{fancyhdr}
% Link "¿Cómo hice este documento?" para el footer-izquierda (solo si hay repo_url).
((* if repo_url *))
\newcommand{\cvfooterlink}{\href{((( repo_url )))}{\footnotesize\color{body!75}\faGithub\, ((( labels.como_hice ))) }}
((* else *))
\newcommand{\cvfooterlink}{}
((* endif *))
% Páginas 2+: marca arriba-derecha + link abajo-izquierda + número abajo-centro.
\fancypagestyle{cvpage}{%
  \fancyhf{}%
  \fancyhead[R]{\footnotesize\color{body!75}Curriculum Vitae - ((( personales.nombre )))}%
  \fancyfoot[L]{\cvfooterlink}%
  \fancyfoot[C]{\footnotesize\color{body!75}\thepage}%
  \renewcommand{\headrulewidth}{0pt}%
  \renewcommand{\footrulewidth}{0pt}%
}
% Primera página: link abajo-izquierda + número abajo-centro (sin marca arriba).
\fancypagestyle{cvfirst}{%
  \fancyhf{}%
  \fancyfoot[L]{\cvfooterlink}%
  \fancyfoot[C]{\footnotesize\color{body!75}\thepage}%
  \renewcommand{\headrulewidth}{0pt}%
  \renewcommand{\footrulewidth}{0pt}%
}

\begin{document}

% === INFORMACIÓN PERSONAL ===
\name{ ((( personales.nombre ))) }
\tagline{ ((( personales.titulo ))) }

% Sin espacios dentro de los {}: altacv los mete en la URL del link y la rompe.
\personalinfo{
  \email{((( personales.correo )))}
  \phone{((( personales.telefono )))}
  ((* if personales.linkedin *))\linkedin{((( personales.linkedin )))}((* endif *))
  ((* if personales.github *))\github{((( personales.github )))}((* endif *))
}

\makecvheader
\pagestyle{cvpage}      % páginas 2+: marca + número
\thispagestyle{cvfirst} % página 1: solo número (sin marca)

% === RESUMEN PROFESIONAL ===
((* if personales.resumen *))
\cvsection{ ((( labels.resumen ))) }
((( personales.resumen )))
\medskip
((* endif *))

% ============================================================================
% BANDA SUPERIOR: Aptitudes Técnicas (70%) | Idiomas (30%)
% ============================================================================
\columnratio{0.7}
\begin{paracol}{2}
\vspace*{0pt}% ancla el tope de la columna: evita que paracol descarte el \bigskip
             % de \cvsection en una columna y no en la otra (líneas desalineadas).

% --- Columna izquierda (70%): Aptitudes (por categorías) ---
% aptitudes.aptitudes es una lista de grupos { categoria, skills }. "categoria"
% es un valor (se traduce); cada grupo = subtítulo + su caja de chips.
\cvsection{ ((( labels.aptitudes ))) }
((* for grupo in aptitudes.aptitudes *))
\cvsubsection{ ((( grupo.categoria ))) }
\cvtagbox{%
((* for apt in grupo.skills *))
\cvtag{ ((( apt ))) }
((* endfor *))
}
((* endfor *))

\switchcolumn
\vspace*{0pt}% mismo ancla en la columna derecha para que ambas líneas alineen

% --- Columna derecha (30%): Idiomas + Intereses (apilados, rellenan el hueco) ---
\cvsection{ ((( labels.idiomas ))) }
\cvtagbox{%
((* for idioma in otros.idiomas *))
\cvtag{ ((( idioma.idioma ))) - ((( idioma.nivel_lectura ))) }
((* endfor *))
}

\medskip

% Intereses: debajo de Idiomas, en la MISMA columna derecha (antes iba a ancho
% completo y dejaba un hueco vacío bajo Idiomas).
\cvsection{ ((( labels.intereses ))) }
\cvtagbox{%
((* for interes in otros.intereses *))
\cvtag{ ((( interes ))) }
((* endfor *))
}

\end{paracol}

\medskip

% ============================================================================
% BLOQUE INFERIOR (ancho completo): Experiencia → Educación → Publicaciones → Certificaciones
% ============================================================================

% === EXPERIENCIA LABORAL ===
\cvsection{ ((( labels.experiencia ))) }
((* for exp in experiencia.experiencia *))
\cvevent{ ((( exp.puesto ))) }{ ((( exp.empresa ))) }{ ((( exp.fecha_inicio ))) -- ((( exp.fecha_fin ))) }{ ((( exp.ubicacion ))) }
((* if exp.descripcion and exp.descripcion|length > 0 *))
\begin{itemize}
  ((* for desc in exp.descripcion *))
  \item ((( desc )))
  ((* endfor *))
\end{itemize}
((* endif *))
((* if not loop.last *))\divider((* endif *))
((* endfor *))

\medskip

% === EDUCACIÓN ===
\cvsection{ ((( labels.educacion ))) }
((* for edu in educacion.educacion *))
\cvevent{ ((( edu.grado ))) }{ ((( edu.institucion ))) }{ ((( edu.fecha_inicio ))) -- ((( edu.fecha_fin ))) }{ ((( edu.ubicacion ))) }
((* if edu.trabajo_grado and edu.trabajo_grado != "Pendiente." *))
((( labels.trabajo_grado ))): ((( edu.trabajo_grado )))
((* endif *))
((* if not loop.last *))\divider((* endif *))
((* endfor *))

\medskip

% === PUBLICACIONES ===
\cvsection{ ((( labels.publicaciones ))) }
((* for pub in publicaciones.publicaciones *))
\cvevent{ ((( pub.titulo ))) }{ ((( pub.autores ))) }{ ((( pub.año ))) }{ ((( pub.tipo ))) }
Revista: ((( pub.revista )))
((* if pub.enlace *))
\href{((( pub.enlace )))}{ ((( labels.ver_pub ))) }
((* endif *))
((* if not loop.last *))\divider((* endif *))
((* endfor *))

\medskip

% === CERTIFICACIONES ===
\cvsection{ ((( labels.certificaciones ))) }
((* for cert in certificaciones.certificaciones *))
\cvevent{ ((( cert.nombre ))) }{ ((( cert.institucion ))) }{ ((( cert.fecha_inicio ))) -- ((( cert.fecha_fin ))) }{}
((* if not loop.last *))\divider((* endif *))
((* endfor *))

\medskip

% === REFERENCIAS (Profesionales | Personales) ===
((* if referencias is defined and (referencias.profesionales or referencias.personales) *))
\needspace{12\baselineskip}% reserva para el título + la banda; si no caben, saltan juntos
\cvsection{ ((( labels.referencias ))) }
\columnratio{0.5}
\begin{paracol}{2}
\vspace*{0pt}

% --- Profesionales ---
((* if referencias.profesionales *))
\cvsubsection{ ((( labels.ref_profesionales ))) }
((* for r in referencias.profesionales *))
\textbf{((( r.nombre )))}\\
((( r.cargo )))((* if r.organizacion *)) · ((( r.organizacion )))((* endif *))\\
{\small ((( r.contacto )))}

\smallskip
((* endfor *))
((* endif *))

\switchcolumn
\vspace*{0pt}

% --- Personales ---
((* if referencias.personales *))
\cvsubsection{ ((( labels.ref_personales ))) }
((* for r in referencias.personales *))
\textbf{((( r.nombre )))}\\
((( r.relacion )))\\
{\small ((( r.contacto )))}

\smallskip
((* endfor *))
((* endif *))

\end{paracol}
((* endif *))

\end{document}